\documentclass[11pt]{article}

\usepackage{amssymb,amsmath,amsfonts,amsthm}
\usepackage{lmodern}
\usepackage{hyperref}
\usepackage{mathtools}
\usepackage{enumerate}
\usepackage{booktabs}
\usepackage{algorithm}
\usepackage{algorithmic}

\setlength{\textwidth}{15.5cm} \setlength{\headheight}{0.5cm} \setlength{\textheight}{21.5cm}
\setlength{\oddsidemargin}{0.25cm} \setlength{\evensidemargin}{0.25cm} \setlength{\topskip}{0.5cm}
\setlength{\footskip}{1.5cm} \setlength{\headsep}{0cm} \setlength{\topmargin}{0.5cm}

\newtheorem{definition}{Definition}[section]
\newtheorem{proposition}{Proposition}[section]
\newtheorem{lemma}{Lemma}[section]
\newtheorem{remark}{Remark}[section]

\newcommand{\R}{\mathbb{R}}
\newcommand{\norm}[1]{\left\lVert#1\right\rVert}
\newcommand{\pbar}{\bar{p}}

\allowdisplaybreaks

\begin{document}

\title{QP-Guided Redistricting via Marked Edge Walks}
\author{Alaittin K\i rt\i \c{s}o\u{g}lu}
\footnotetext[1]{Department of Applied Mathematics, Illinois Institute of Technology}
\maketitle


%═══════════════════════════════════════════════════════════════════════════════
\section{Setup}\label{sec:setup}
%═══════════════════════════════════════════════════════════════════════════════

Let $G = (V,E)$ be a connected graph, $|V| = n$. Each vertex $v$ has population $p_v > 0$. Fix $k$ districts with ideal population $\pbar = \frac{1}{k}\sum_v p_v$ and balance tolerance $\varepsilon > 0$.

A \textbf{partition state} is a pair $(T, M)$ where:
\begin{itemize}
\item $T$ is a spanning tree of $G$.
\item $M \subset E(T)$ with $|M| = k-1$ is the \textbf{marked set}.
\end{itemize}
Removing $M$ from $T$ yields $k$ connected subtrees --- the \textbf{districts} $D_1, \ldots, D_k$. Since $T$ is a tree, every choice of $k-1$ edges gives a valid partition; contiguity is automatic.

The \textbf{feasible set} is
\[
\mathcal{B} = \bigl\{(T, M) : |p(D_i) - \pbar| \leq \varepsilon\, \pbar \text{ for all } i\bigr\}.
\]

We encode $M$ as a binary vector $x \in \{0,1\}^{|E(T)|}$ with $x_e = 1$ iff $e \in M$, and relax to
\[
\mathcal{X} = \bigl\{x \in [0,1]^{|E(T)|} : \mathbf{1}^\top x = k-1\bigr\}.
\]


%═══════════════════════════════════════════════════════════════════════════════
\section{Subtree Populations and the Balance Signal}\label{sec:subtree}
%═══════════════════════════════════════════════════════════════════════════════

Root the spanning tree $T$ at an arbitrary vertex $r$. For each non-root vertex $v$, let $e_v = (v, \mathrm{parent}(v))$ denote the edge connecting $v$ to its parent, and define the \textbf{subtree population}
\[
\sigma_v \;=\; \sum_{u \in \mathrm{sub}(v)} p_u,
\]
where $\mathrm{sub}(v)$ is the set of all descendants of $v$ (including $v$ itself). This quantity is computable once from $T$ and $\mathbf{p}$ in $O(n)$ time via a bottom-up pass, and is independent of $M$.

\begin{definition}[Balance signal]\label{def:balance-signal}
The \textbf{balance signal} is the diagonal matrix
\[
\Sigma = \mathrm{diag}\bigl(\sigma_v - \pbar\bigr)_{v \neq r} \;\in\; \R^{(n-1) \times (n-1)}.
\]
\end{definition}

When edge $e_v$ is cut (i.e., $x_{e_v} = 1$) and no ancestor edge is also cut, the component below $v$ has population exactly $\sigma_v$. The deviation from $\pbar$ is $\sigma_v - \pbar$. This motivates the \textbf{linearized balance 1-form}:
\[
\omega(\mathbf{x}) = \Sigma \mathbf{x}.
\]
Perfect balance holds if and only if $\omega(\mathbf{x}^*) = \mathbf{0}$, i.e., every cut subtree has population exactly $\pbar$.

\begin{remark}[Root dependence]
The subtree populations $\sigma_v$ depend on the choice of root $r$. Different roots assign different $\sigma_v$ to the same edge. We randomize the root at each step, which adds stochasticity to the QP and prevents the proposal from systematically favoring edges near one end of the tree.
\end{remark}


%═══════════════════════════════════════════════════════════════════════════════
\section{The Edge Laplacian}\label{sec:edge-laplacian}
%═══════════════════════════════════════════════════════════════════════════════

Let $B_1 \in \R^{n \times (n-1)}$ be the signed node-edge incidence matrix of $T$: column $e_v$ has $+1$ at row $v$ and $-1$ at row $\mathrm{parent}(v)$. The \textbf{edge Laplacian} (Hodge 1-Laplacian on the tree) is
\[
\Delta_1^T = B_1^\top B_1 \;\in\; \R^{(n-1) \times (n-1)}.
\]
Since $T$ has no cycles, $\Delta_1^T$ reduces to the combinatorial edge Laplacian. Its diagonal entries are $(\Delta_1^T)_{ee} = \deg_T(v)$, the degree of vertex $v$ in the tree, and its off-diagonal entries encode adjacency between edges sharing a vertex.

The edge Laplacian enters the QP through the \textbf{balance energy}:
\[
\norm{\omega(\mathbf{x})}^2_{\Delta_1} \;=\; \mathbf{x}^\top \bigl(\Sigma\, \Delta_1^T\, \Sigma\bigr)\, \mathbf{x}.
\]
This measures how much the balance signal varies across the tree structure. It penalizes marked sets where adjacent tree edges have very different subtree population deviations.


%═══════════════════════════════════════════════════════════════════════════════
\section{Degeneracy Factor}\label{sec:degeneracy}
%═══════════════════════════════════════════════════════════════════════════════

Multiple lifted states $(T, M) \in \mathcal{X}$ can correspond to the same partition $\xi \in \mathcal{P}_k(G)$. To sample from a desired distribution $\pi(\xi)$ over partitions, we must account for this multiplicity.

\begin{definition}[Degeneracy factor]\label{def:degeneracy}
For a connected $k$-partition $\xi = (D_1, \ldots, D_k)$, the \textbf{degeneracy factor} is
\[
\tau(\xi) \;=\; \prod_{i=1}^{k} t(D_i) \;\cdot\; \prod_{i \sim j} |\partial(D_i, D_j)|,
\]
where:
\begin{itemize}
\item $t(D_i)$ is the number of spanning trees of the subgraph $G[D_i]$ induced by district $D_i$,
\item $|\partial(D_i, D_j)|$ is the number of edges in $G$ between adjacent districts $D_i$ and $D_j$.
\end{itemize}
\end{definition}

The factor $\tau(\xi)$ counts the number of states $(T, M)$ in $\mathcal{X}$ that map to the same partition $\xi$. To see this: any spanning tree $T$ compatible with $\xi$ must (i)~restrict to a spanning tree within each district (giving $\prod_i t(D_i)$ choices), and (ii)~connect adjacent districts by exactly one edge each (giving $\prod_{i \sim j} |\partial(D_i, D_j)|$ choices for the marked edges).

\begin{remark}[Why $\tau$ matters]
Without the $\tau$ correction, MEW samples from the \emph{lifted} distribution on $\mathcal{X}$, which heavily oversamples partitions with many spanning trees. On real dual graphs this bias is extreme: for Texas congressional districts, the degeneracy factors of two plans can differ by a factor of $10^{466}$ \cite{mew}. The target distribution on $\mathcal{X}$ is chosen as
\[
p(T, M) \;\propto\; \frac{\exp(J(\xi(T,M)))}{\tau(\xi(T,M))},
\]
so that the induced distribution on partitions is $\pi(\xi) \propto \exp(J(\xi))$, as desired. Dividing by $\tau$ exactly cancels the spanning-tree multiplicity.
\end{remark}

\begin{remark}[Computation]
The spanning tree count $t(D_i)$ is computed via the Matrix-Tree Theorem (Kirchhoff's theorem): $t(D_i) = \det(L_i^*)$, where $L_i^*$ is any cofactor of the graph Laplacian $L_i$ of $G[D_i]$. In practice, we compute $\log t(D_i) = \sum_j \log \lambda_j$, where $\lambda_j$ are the nonzero eigenvalues of $L_i$, or equivalently via LU decomposition of $L_i^*$. The cost is $O(|D_i|^3)$ per district, giving $O(n_{\max}^3)$ overall, where $n_{\max} = \max_i |D_i|$. This is the dominant cost in both MEW and QP-MEW.
\end{remark}

\begin{remark}[Interaction with the QP]
The degeneracy correction appears in the MH acceptance ratio as $\log \tau(\xi) - \log \tau(\xi')$. If the QP proposal jumps far from the current partition, the new partition may have a very different $\tau$, causing the ratio to be large and the proposal to be rejected. The locality term $\lambda \Sigma^2$ in the QP limits step size, keeping $\tau(\xi') \approx \tau(\xi)$ and maintaining reasonable acceptance rates.
\end{remark}


%═══════════════════════════════════════════════════════════════════════════════
\section{The QP Proposal}\label{sec:qp}
%═══════════════════════════════════════════════════════════════════════════════

After a cycle basis step produces a new tree $T'$, we solve a QP over $\mathbf{x} \in \mathcal{X}$ to select the new marked set $M'$. The QP has two terms: a balance energy that steers toward feasible partitions, and a locality term that keeps $M'$ close to the current $M$.

\subsection{The full QP}

\begin{equation}\label{eq:qp}
\boxed{\min_{\mathbf{x} \in \mathcal{X}} \;\; \mathbf{x}^\top Q\, \mathbf{x} \;-\; 2\lambda\, \mathbf{q}^\top \mathbf{x}}
\end{equation}
where
\begin{align}
Q &= \Sigma\, \Delta_1^T\, \Sigma \;+\; \lambda\, \Sigma^2, \label{eq:Q} \\
\mathbf{q} &= \Sigma^2\, \mathbf{x}_M, \label{eq:q}
\end{align}
and $\mathbf{x}_M \in \{0,1\}^{n-1}$ encodes the current marked set.

\medskip
\noindent\textbf{Interpretation of each term:}
\begin{itemize}
\item \textbf{Balance energy} $\mathbf{x}^\top (\Sigma\, \Delta_1^T\, \Sigma)\, \mathbf{x}$: penalizes marked sets that create large population imbalances. At the minimizer, every cut subtree has population close to $\pbar$.
\item \textbf{Locality} $\lambda\, \mathbf{x}^\top \Sigma^2\, \mathbf{x} - 2\lambda\, \mathbf{q}^\top \mathbf{x} = \lambda\, (\mathbf{x} - \mathbf{x}_M)^\top \Sigma^2\, (\mathbf{x} - \mathbf{x}_M) - \lambda\, \mathbf{x}_M^\top \Sigma^2\, \mathbf{x}_M$: keeps $M'$ close to $M$ in the population geometry. The $\Sigma^2$ weighting means that edges with large subtree population deviation are expensive to swap.
\item \textbf{$\lambda$}: controls the balance-locality tradeoff. $\lambda = 0$ gives the globally most balanced partition on $T'$. $\lambda \to \infty$ gives $M' = M$ (no move). Finite $\lambda$ gives a local, approximately balanced perturbation of $M$.
\end{itemize}

\subsection{Greedy solution}

For binary $\mathbf{x}$ with $\mathbf{1}^\top \mathbf{x} = k-1$, the QP reduces to selecting the $k-1$ edges with smallest marginal cost. The diagonal approximation gives:
\begin{equation}\label{eq:score}
\mathrm{score}(e_v) \;=\; (\sigma_v - \pbar)^2 \,\bigl(\deg_T(v) \;+\; \lambda\,(1 - 2\,(\mathbf{x}_M)_{e_v})\bigr).
\end{equation}
\begin{itemize}
\item If $e_v \in M$ (currently marked): score $= (\sigma_v - \pbar)^2\,(\deg_T(v) - \lambda)$. The $-\lambda$ term reduces the score, favoring retention.
\item If $e_v \notin M$ (not marked): score $= (\sigma_v - \pbar)^2\,(\deg_T(v) + \lambda)$. The $+\lambda$ term increases the score, discouraging addition.
\item Edges with $\sigma_v \approx \pbar$ have near-zero score regardless of $\lambda$: perfectly balanced cuts are always preferred.
\end{itemize}

The greedy solve picks the $k-1$ edges with smallest score. This is $O(n \log n)$ (sorting) after an $O(n)$ preprocessing step (BFS rooting + bottom-up subtree population computation).

\subsection{Stochasticity}

Two sources of stochasticity ensure the proposal kernel has full support over $\mathcal{F}_k(T') = \{M' \subseteq E(T') : |M'| = k-1\}$:

\begin{enumerate}
\item \textbf{$\lambda$-perturbation.} At each step, $\lambda$ is sampled from a Gamma distribution:
\[
\lambda \sim \mathrm{Gamma}\!\left(\frac{1}{\nu^2},\; \frac{\lambda_0 \nu^2}{1}\right), \qquad \mathbb{E}[\lambda] = \lambda_0, \quad \mathrm{CV}[\lambda] = \nu.
\]
Small $\lambda$ draws produce exploratory moves (jump toward the global balance optimum); large $\lambda$ draws produce conservative moves (stay near $M$).

\item \textbf{Random root.} The root $r$ is chosen uniformly at random at each step. Since $\sigma_v$ depends on the root, different roots produce different QP scores and hence different $M'$. This adds a natural, zero-cost source of diversity.
\end{enumerate}


%═══════════════════════════════════════════════════════════════════════════════
\section{The Full Algorithm}\label{sec:algorithm}
%═══════════════════════════════════════════════════════════════════════════════

\begin{algorithm}[H]
\caption{QP-Guided MEW Step}
\label{alg:qp-mew}
\begin{algorithmic}[1]
\REQUIRE State $(T, M)$, graph $G$, parameters $\lambda_0, \nu, \varepsilon, k$
\ENSURE New state $(T', M')$ or rejection (stay at $(T, M)$)

\medskip
\STATE \textbf{Cycle basis step} (identical to plain MEW):
\STATE \quad Choose $e^+ \in E(G) \setminus E(T)$ uniformly at random
\STATE \quad Find the unique cycle $C$ in $T \cup \{e^+\}$
\STATE \quad Choose $e^- \in C \setminus M$ uniformly at random
\STATE \quad $T' \leftarrow (T \cup \{e^+\}) \setminus \{e^-\}$

\medskip
\STATE \textbf{QP-guided marked edge selection:}
\STATE \quad Choose root $r$ uniformly at random from $V$
\STATE \quad Compute subtree populations $\sigma_v$ for all $v \neq r$ via bottom-up pass on $T'$
\STATE \quad Sample $\lambda \sim \mathrm{Gamma}(1/\nu^2, \lambda_0 \nu^2)$
\STATE \quad Compute $\mathrm{score}(e_v) = (\sigma_v - \pbar)^2\,(\deg_{T'}(v) + \lambda(1 - 2\,(\mathbf{x}_M)_{e_v}))$ for each $v \neq r$
\STATE \quad $M' \leftarrow$ the $k-1$ edges with smallest score

\medskip
\STATE \textbf{Metropolis--Hastings acceptance:}
\STATE \quad Compute partition $\xi' = T' \setminus M'$
\STATE \quad \textbf{if} any district $D_i$ has $|p(D_i) - \pbar| > \varepsilon\, \pbar$ \textbf{then reject}
\STATE \quad Compute transition ratio: $\log r_{\mathrm{cycle}} = \log|C \setminus M| - \log|C' \setminus M'|$
\STATE \quad Compute energy: $J(\xi) = J_{\mathrm{balance}}(\xi) + J_{\mathrm{compact}}(\xi)$
\STATE \quad Compute degeneracy: $\tau(\xi) = \prod_i t(\xi_i) \cdot \prod_{i \sim j} |\partial(\xi_i, \xi_j)|$
\STATE \quad $\log \alpha = (J(\xi') - J(\xi)) + (\log \tau(\xi) - \log \tau(\xi')) + \log r_{\mathrm{cycle}}$
\STATE \quad Accept with probability $\min(1, e^{\alpha})$
\end{algorithmic}
\end{algorithm}


%═══════════════════════════════════════════════════════════════════════════════
\section{Energy Function}\label{sec:energy}
%═══════════════════════════════════════════════════════════════════════════════

The MH acceptance uses an energy function $J(\xi) = J_{\mathrm{balance}}(\xi) + J_{\mathrm{compact}}(\xi)$ that serves two purposes: the balance term keeps the chain inside $\mathcal{B}$, and the compactness term provides a gradient for exploration.

\subsection{Balance term: soft barrier}

We use a flat-bottom potential that is zero inside $\mathcal{B}$ and turns on quadratically outside:
\[
J_{\mathrm{balance}}(\xi) \;=\; -\beta_b \sum_{i=1}^{k} \max\!\left(0,\; \frac{|p(D_i) - \pbar|}{\varepsilon\,\pbar} - 1\right)^{\!2}.
\]
This is identically zero when all districts satisfy the $\varepsilon$-tolerance. Outside $\mathcal{B}$, it provides a strong restoring force ($\beta_b \gg 1$). Unlike a quadratic penalty centered at $\pbar$, this does not bias the chain toward perfectly balanced partitions --- all states in $\mathcal{B}$ are treated equally.

\subsection{Compactness term: exploration gradient}

Without a compactness gradient, the chain has no preferred direction inside $\mathcal{B}$ and mixing is slow. We add
\[
J_{\mathrm{compact}}(\xi) \;=\; -\beta_c\, (c(\xi) - \mu_c)^2,
\]
where $c(\xi) = |\{(u,v) \in E(G) : \xi(u) \neq \xi(v)\}|$ is the number of cut edges and $\mu_c$ is a target cut-edge count. This gives the chain a landscape to explore along: it moves toward partitions with approximately $\mu_c$ cut edges, providing directionality that accelerates mixing.


%═══════════════════════════════════════════════════════════════════════════════
\section{Computational Cost}\label{sec:cost}
%═══════════════════════════════════════════════════════════════════════════════

The per-step cost of QP-guided MEW consists of:
\begin{center}
\renewcommand{\arraystretch}{1.3}
\begin{tabular}{lll}
\toprule
\textbf{Operation} & \textbf{Cost} & \textbf{Shared with MEW?} \\
\midrule
Cycle basis step & $O(n)$ & Yes \\
BFS rooting of $T'$ & $O(n)$ & No (QP only) \\
Subtree populations & $O(n)$ & No (QP only) \\
Edge scoring + sort & $O(n \log n)$ & No (QP only) \\
Partition extraction & $O(n)$ & Yes \\
Population check & $O(n)$ & Yes \\
Degeneracy $\tau(\xi)$ & $O(n_{\max}^3)$ & Yes \\
\bottomrule
\end{tabular}
\end{center}
where $n_{\max}$ is the size of the largest district. The QP overhead (rooting + subtree populations + scoring) is $O(n \log n)$, which is dominated by the $O(n_{\max}^3)$ degeneracy computation that both MEW and QP-MEW must perform. The degeneracy computation is the true bottleneck for large graphs.


%═══════════════════════════════════════════════════════════════════════════════
\section{Relation to the Population Geometry}\label{sec:geometry}
%═══════════════════════════════════════════════════════════════════════════════

The QP has a natural geometric interpretation. Define the \textbf{population deviation signal} for a marked set $M$ as
\[
\omega_M \;=\; \Sigma\, \mathbf{x}_M \;\in\; \R^{n-1}.
\]
The $e$-th component $(\omega_M)_e = (\sigma_v - \pbar) \cdot (\mathbf{x}_M)_e$ is the deviation from $\pbar$ of the subtree below $v$, weighted by whether edge $e$ is cut.

The QP objective~\eqref{eq:qp} can be written as
\[
\norm{\omega_{M'}}^2_{\Delta_1} \;+\; \lambda\, \norm{\omega_{M'} - \omega_M}^2,
\]
where $\norm{\cdot}_{\Delta_1}$ is the norm induced by the edge Laplacian. The first term minimizes the balance energy of $M'$; the second keeps $M'$ close to $M$ in the population geometry.

\begin{definition}[Population distance]\label{def:distance}
For two marked sets $M, M'$ on the same tree $T$, the \textbf{population distance} is
\[
d(M, M') = \norm{\omega_M - \omega_{M'}}.
\]
\end{definition}

This distance measures how much the population balance changes when moving from $M$ to $M'$. Two marked sets that differ in many edges but produce similar population splits are close; a single-edge swap that creates a large imbalance is far. The QP proposals move in directions of small population distance.


%═══════════════════════════════════════════════════════════════════════════════
\section{Properties}\label{sec:properties}
%═══════════════════════════════════════════════════════════════════════════════

\begin{proposition}[Convexity]\label{prop:convex}
The matrix $Q = \Sigma\,\Delta_1^T\,\Sigma + \lambda\,\Sigma^2$ is positive semidefinite for all $\lambda \geq 0$, so the QP~\eqref{eq:qp} is convex.
\end{proposition}

\begin{proof}
$\Delta_1^T = B_1^\top B_1 \succeq 0$ and $\Sigma^2 \succeq 0$, so $\Sigma\,\Delta_1^T\,\Sigma = \Sigma B_1^\top B_1 \Sigma = (B_1 \Sigma)^\top (B_1 \Sigma) \succeq 0$, and $Q$ is the sum of two PSD matrices.
\end{proof}

\begin{proposition}[Behavior at extremes]\label{prop:extremes}
\begin{enumerate}[(i)]
\item At $\lambda = 0$: the minimizer selects the most balanced partition of $T'$ (globally optimal cut).
\item As $\lambda \to \infty$: the minimizer converges to $\mathbf{x}_M$ (no move).
\item At finite $\lambda$: the minimizer is a locally balanced perturbation of $M$ --- exactly the type of proposal MEW needs to reduce rejection.
\end{enumerate}
\end{proposition}

\begin{proposition}[Feasibility]\label{prop:feasibility}
For any $\mathbf{x}^* \in \mathcal{X}$, the rounded solution $\hat{\mathbf{x}}$ (top $k-1$ entries set to 1, rest to 0) satisfies $|\hat{M}| = k-1$, so $T' \setminus \hat{M}$ is always a valid $k$-partition with $k$ connected districts. No additional connectivity check is needed.
\end{proposition}


%═══════════════════════════════════════════════════════════════════════════════
\section{Discussion}\label{sec:discussion}
%═══════════════════════════════════════════════════════════════════════════════

\noindent\textbf{Why MEW rejects.} Plain MEW proposes marked edge changes uniformly at random. On a tree with $n - 1$ edges, most choices of $k - 1$ edges produce imbalanced partitions. The rejection rate grows with $n$ because the feasible set $\mathcal{B}$ becomes a smaller fraction of the state space.

\medskip
\noindent\textbf{What the QP does.} The QP replaces uniform selection with informed selection. By scoring each edge using the subtree population $\sigma_v$, it identifies edges that create balanced cuts. The locality term ensures the proposal stays close to the current state, so the MH degeneracy correction $\tau(\xi)/\tau(\xi')$ remains close to 1.

\medskip
\noindent\textbf{The $\lambda$ tradeoff.} Small $\lambda$ produces balanced but distant proposals (large $\tau$ ratio, low acceptance). Large $\lambda$ produces local but potentially imbalanced proposals (small $\tau$ ratio, high acceptance but slow mixing). The optimal $\lambda$ balances these effects, maximizing the effective sample size per unit time.

\medskip
\noindent\textbf{Comparison with plain MEW.}
\begin{center}
\renewcommand{\arraystretch}{1.3}
\begin{tabular}{lll}
\toprule
& \textbf{MEW} & \textbf{QP-MEW} \\
\midrule
Tree step & Cycle basis (uniform) & Cycle basis (uniform) \\
Marked edge step & Random slide & QP-guided selection \\
Balance awareness & None (hard rejection) & Built into QP via $\Sigma$ \\
Stochasticity & Uniform random & $\lambda$-perturbation $+$ random root \\
Per-step cost & $O(n)$ & $O(n \log n)$ \\
\bottomrule
\end{tabular}
\end{center}


\bibliographystyle{plain}
\bibliography{ref}

\end{document}
